\section{延伸分析：哪些判断最值得继续验证}

\subsection{验证点一：中国 AI 需求到底像 SaaS 还是 Cloud}

如果 AI 支出像 SaaS，中国市场可能受限于传统软件付费文化；如果像 cloud，它可能成为研发、客服、数据分析、搜索推荐和自动化的基础设施。可以观察：

\begin{itemize}[leftmargin=1.5em]
    \item 云厂商模型调用收入是否增长；
    \item 企业私有化部署和行业大模型项目是否持续付费；
    \item coding agent、客服 agent、办公 agent 是否进入高频工作流；
    \item 推理芯片和模型服务是否形成规模化价格竞争。
\end{itemize}

\subsection{验证点二：开放权重能否形成全球开发者 moat}

中国模型开放权重越强，全球开发者越可能在本地部署、微调和二次开发中选择它们。可观察：

\begin{itemize}[leftmargin=1.5em]
    \item Hugging Face、GitHub、OpenRouter、云市场中的使用量；
    \item vLLM、SGLang、llama.cpp、Ollama 等推理生态支持速度；
    \item 企业是否把中国开放模型作为私有部署默认候选；
    \item 美国政策是否限制开放模型发布，反而给中国模型生态留出空间。
\end{itemize}

\subsection{验证点三：中国能否从 fast follower 变成范式定义者}

这是全文最大开放问题。可以看以下信号：

\begin{itemize}[leftmargin=1.5em]
    \item 是否提出新的训练范式，而不只是优化既有路线；
    \item 是否产生国际社区主动跟随的 agent / RL / multimodal 方法；
    \item 是否在真实产品反馈上形成独特数据优势；
    \item 是否有足够基础研究空间，而不是所有人才都被短期模型竞赛吸走；
    \item 是否能在 Nvidia 受限下发展出不同的高效训练和推理技术路线。
\end{itemize}

\begin{importantbox}{从追赶到定义}
fast-following 能让中国模型持续接近前沿，但范式定义需要另一套激励：允许失败、奖励原创问题、保留基础研究人才，并让研究者不只优化已经被证明有效的路线。
\end{importantbox}

\subsection{验证点四：美国开放模型战略会如何变化}

作者最担心的是，美国一方面希望保持 AI 领导力，另一方面又可能通过行政命令或监管限制开放模型。这可能导致一个悖论：美国闭源模型继续领先，但全球开放开发者生态逐渐被中国模型占据。

\begin{warningbox}{领导力不只等于最强闭源模型}
AI 领导力包括模型能力、开发者生态、标准、工具链、部署框架、硬件优化、教育资源和全球信任。如果美国只保留闭源能力，而放弃开放生态领导力，它可能失去一部分长期影响力。
\end{warningbox}

\subsection{本章小结}

这篇文章最适合转化为四个观察题：AI 需求像什么、开放权重能否形成 moat、中国能否定义新范式、美国是否愿意领导开放生态。


